\begin{textsample}{POS Dim 1 – vem\_ed\_30 – Score 60.00 – vem\_ed\_30/t159.txt}  \label{ex:f1_pos_025}
IA nos Projetos Colaborativos Internacionais (PCIs) Observamos um grande entusiasmo com a utilização das Inteligências Artificiais Generativas (IAGen) e seu potencial de \textbf{expansão} do trabalho individual.
Essas tecnologias estão redefinindo as possibilidades educacionais, oferecendo ferramentas \textbf{que} podem personalizar o \textbf{aprendizado} e amplificar as \textbf{capacidades} tanto de educadores \textbf{quanto} de estudantes.
Paralelamente a esse entusiasmo, existe um medo real sobre a perda de empregos em \textbf{todos} os \textbf{níveis}, inclusive dos próprios programadores.
Esta preocupação \textbf{é} legítima e demanda uma \textbf{reflexão} cuidadosa sobre como preparar nossa força de trabalho para um \textbf{futuro} em constante \textbf{transformação} tecnológica.
Para o Centro Paula Souza, um dos aspectos mais importantes envolve a inclusão proporcionada por essas ferramentas e a maneira de incorporá-las nas práticas educacionais.
O objetivo \textbf{é} garantir \textbf{que} os alunos cheguem mais preparados para o \textbf{mercado} de trabalho, o \textbf{que} faz \textbf{parte} do DNA vocacional do Centro Paula Souza.
A experiência de 12 anos trabalhando com Projetos Colaborativos Internacionais (PCIs) nas Fatecs sempre \textbf{foi} positivamente surpreendente.
Essa trajetória \textbf{consolidada} nos \textbf{permite} \textbf{observar} com clareza tanto os avanços \textbf{quanto} os desafios \textbf{que} enfrentamos na colaboração internacional. \textbf{continuação} Os \textbf{impactos} dos PCIs na formação dos estudantes \textbf{são} multidimensionais.
Há o \textbf{desenvolvimento} de uma \textbf{perspectiva} global sobre \textbf{questões} \textbf{locais}, \textbf{permitindo} \textbf{que} compreendam como problemas \textbf{ambientais}, por exemplo, \textbf{são} universais, mas requerem soluções adaptadas aos \textbf{contextos} \textbf{específicos}.
Além disso, a experiência de trabalhar com equipes internacionais desenvolve \textbf{competências} \textbf{essenciais} para o \textbf{mercado} de trabalho globalizado, incluindo \textbf{flexibilidade} cultural, comunicação eficaz em \textbf{ambientes} diversos e \textbf{capacidade} de liderança em \textbf{contextos} multiculturais.
A promoção da cidadania global por meio dos PCIs \textbf{é} evidente na \textbf{forma} como os estudantes passam a compreender sua responsabilidade não apenas como \textbf{profissionais}, mas como cidadãos do mundo, \textbf{conectando} \textbf{conhecimento} técnico com propósito social e \textbf{impacto} comunitário.
Sempre \textbf{contamos} com o \textbf{apoio} dos professores de Inglês e de Espanhol, \textbf{que} facilitam as mediações de professores e alunos das Fatecs em língua estrangeira, \textbf{criando} \textbf{pontes} \textbf{essenciais} para a comunicação intercultural.
Contudo, uma espécie de “complexo de vira-lata"quase sempre pairava nas \textbf{interações} – uma humildade desnecessária frente a quem \textbf{utiliza} o espanhol ou o inglês diariamente.
Este fenômeno cultural brasileiro muitas vezes limitava o potencial pleno de nossos estudantes, \textbf{criando} barreiras psicológicas \textbf{que} iam além das \textbf{questões} puramente \textbf{linguísticas}.
O \textbf{apoio} dos professores de \textbf{idiomas} e a IAGen têm servido como \textbf{elemento} facilitador para alunos e professores.
A IAGen oferece uma democratização \textbf{linguística} \textbf{significativa}, \textbf{reduzindo} as barreiras de comunicação e aumentando a confiança dos estudantes para participar ativamente das colaborações internacionais.
Observamos uma melhoria notável na qualidade da comunicação, com maior precisão e clareza nas \textbf{interações}, além de uma aceleração considerável nos \textbf{processos} \textbf{colaborativos}, otimizando o tempo dedicado à colaboração efetiva. \textbf{continuação} Essa \textbf{transformação} tecnológica também demanda uma \textbf{mudança} fundamental no paradigma avaliativo.
A verificação do trabalho entregue pelos alunos tem \textbf{que} \textbf{ser} mudada também.
Os professores precisam conhecer o \textbf{que} a IAGen oferece e acompanhar o trabalho passo a passo e não somente no \textbf{momento} da entrega.
Isso requer o \textbf{desenvolvimento} de novas \textbf{competências} de supervisão e orientação, bem como uma \textbf{compreensão} profunda das \textbf{capacidades} e limitações das ferramentas de IA disponíveis.
Do \textbf{ponto} de vista de \textbf{coordenação} de PCIs, \textbf{que} estão inseridos na área de Apoio à Internacionalização do Ensino Superior da Divisão de Extensão e Pesquisa no Ensino Superior (Depes) da Coordenadoria Geral de Ensino Superior de Graduação (CGESG) do CPS, o objetivo principal da equipe \textbf{é} facilitar o trabalho do professor com a IAGen em aspectos como proposição de ideias, resumos de reuniões e escrita de Plano de Trabalho para o PCI.
Esta \textbf{estratégia} visa \textbf{permitir} \textbf{que} o professor da Fatec tenha mais tempo para conversar, conhecer e formar laços mais sólidos com os parceiros internacionais - o \textbf{verdadeiro} coração das colaborações internacionais.
Os principais desafios \textbf{que} enfrentamos incluem a \textbf{necessidade} de capacitação docente para uso efetivo da IAGen, a manutenção da autenticidade nas \textbf{interações} \textbf{humanas}, o \textbf{desenvolvimento} de novos métodos avaliativos e o equilíbrio delicado entre tecnologia e relações \textbf{humanas}.
Por outro lado, as oportunidades emergentes \textbf{são} promissoras: \textbf{expansão} para novos países parceiros, \textbf{desenvolvimento} de projetos mais ambiciosos e complexos, maior inclusão de estudantes com diferentes perfis \textbf{linguísticos} e fortalecimento da presença internacional das Fatecs.
A IAGen não substitui a essência \textbf{humana} da colaboração internacional, mas potencializa nossa \textbf{capacidade} de \textbf{conectar}, \textbf{criar} e transformar por meio da educação globalizada.
A experiência demonstra \textbf{que}, quando bem \textbf{implementada}, a tecnologia serve como uma \textbf{ponte} \textbf{que} aproxima culturas e democratiza oportunidades educacionais, \textbf{permitindo} \textbf{que} os estudantes participem de \textbf{forma} mais plena e confiante no cenário acadêmico internacional.

% matched lemmas: ambiental, ambiente, apoio, aprendizado, capacidade, colaborativo, competência, compreensão, conectar, conhecimento, consolidar, contar, contexto, continuação, coordenação, criar, desenvolvimento, elemento, específico, essencial, estratégia, expansão, flexibilidade, forma, futuro, humano, idioma, impacto, implementar, interação, linguístico, local, mercado, momento, mudança, necessidade, nível, observar, parte, permitir, perspectiva, ponte, ponto, processo, profissional, quanto, que, questão, reduzir, reflexão, ser, significativo, todo, transformação, utilizar, verdadeiro
\end{textsample}
